% main.tex — SF-Harris Report: Main document
% Compile with: pdflatex main && bibtex main && pdflatex main && pdflatex main

\documentclass[12pt,letterpaper]{report}
% preamble.tex — Packages, macros, and title page for SF-Harris report

% ==============================================================================
% Packages
% ==============================================================================
\usepackage[utf8]{inputenc}
\usepackage[T1]{fontenc}
\usepackage[spanish]{babel}
\usepackage{amsmath,amssymb,amsthm}
\usepackage{mathtools}
\usepackage{booktabs}
\usepackage{graphicx}
\usepackage{hyperref}
\usepackage{cleveref}
\usepackage{enumitem}
\usepackage{float}
\usepackage{subcaption}
\usepackage{geometry}
\usepackage{fancyhdr}
\usepackage{natbib}
\usepackage{xcolor}

\geometry{letterpaper, margin=1in, headheight=14.5pt}

% Only show chapters in TOC (no sections)
\setcounter{tocdepth}{0}

% ==============================================================================
% Theorem environments (Spanish names)
% ==============================================================================
\newtheorem{teorema}{Teorema}
\newtheorem{proposicion}[teorema]{Proposici\'{o}n}
\newtheorem{lema}[teorema]{Lema}
\newtheorem{corolario}[teorema]{Corolario}
\newtheorem{definicion}{Definici\'{o}n}
\newtheorem{observacion}{Observaci\'{o}n}

% ==============================================================================
% Macros
% ==============================================================================
\newcommand{\E}{\mathbb{E}}
\newcommand{\Var}{\operatorname{Var}}
\newcommand{\Prob}{\mathbb{P}}
\newcommand{\R}{\mathbb{R}}
\newcommand{\N}{\mathcal{N}}
\newcommand{\iid}{\overset{\text{iid}}{\sim}}
\newcommand{\pstay}{p_{\text{stay}}}
\newcommand{\AAD}{\operatorname{AAD}}
\newcommand{\RV}{\operatorname{RV}}
\newcommand{\GIG}{\operatorname{GIG}}

% ==============================================================================
% Hyperref setup
% ==============================================================================
\hypersetup{
    colorlinks=true,
    linkcolor=blue!70!black,
    citecolor=green!50!black,
    urlcolor=blue!70!black,
}

% ==============================================================================
% Title page
% ==============================================================================
\title{
    \vspace{-2cm}
    \textbf{El Modelo SF-Harris: \\
    Replicaci\'{o}n, An\'{a}lisis Cr\'{\'i}tico \\
    y Reducci\'{o}n al Modelo Esencial}
}
\author{[Nombre del Autor] \\ \smallskip \small Universidad [Nombre] \\ Departamento de [Departamento]}
\date{[Fecha]}

% ==============================================================================
% Headers
% ==============================================================================
\pagestyle{fancy}
\fancyhf{}
\fancyhead[L]{\leftmark}
\fancyhead[R]{\thepage}
\renewcommand{\headrulewidth}{0.4pt}

\begin{document}

\maketitle
\thispagestyle{empty}

\begin{abstract}
Este trabajo replica y extiende el modelo de volatilidad estoc\'{a}stica SF-Harris propuesto por Anzarut (2024).
Replicamos la Tabla~3 de la tesis original, logrando una desviaci\'{o}n absoluta promedio (AAD) de 0.3\,pp con barras de d\'{o}lar,
superando los 0.8\,pp reportados originalmente.
A trav\'{e}s de once pruebas experimentales, demostramos que la cadena de Harris es \textbf{decorativa}:
no aporta poder predictivo m\'{a}s all\'{a} del que ya proporciona la distribuci\'{o}n emp\'{\'i}rica de los datos de entrenamiento.
Identificamos dos ingredientes esenciales: (1)~muestreo emp\'{\'i}rico de $Q$ y (2)~estimaci\'{o}n bayesiana de par\'{a}metros $(\mu, \sigma)$
para centrar y escalar la distribuci\'{o}n predictiva.
El par\'{a}metro $\pstay$, que la literatura se\~{n}ala como crucial, resulta no informativo:
su variaci\'{o}n de $0.0$ a $0.95$ dentro del muestreador de Gibbs produce AAD entre $0.41$ y $0.54$\,pp (esencialmente plano).
Discutimos aplicaciones financieras en volatilidad de volatilidad, swaps de varianza y opciones sobre VIX.
\end{abstract}

\tableofcontents
\newpage

% 01-introduccion.tex — Chapter 1: Introducci\'{o}n

\chapter{Introducci\'{o}n}

\section{Motivaci\'{o}n}

La volatilidad estoc\'{a}stica es un componente central de los modelos financieros modernos.
Desde la opci\'{o}n de Black-Scholes hasta los modelos Heston y GARCH, la volatilidad
se modela como un proceso latente que evoluciona en el tiempo.
Sin embargo, la \emph{predicci\'{o}n de cobertura} ---qu\'{e} tan bien un modelo calibra
intervalos de probabilidad para la volatilidad futura--- recibe menos atenci\'{o}n que
la predicci\'{o}n puntual.

El modelo SF-Harris (volatilidad estoc\'{a}stica con recurrencia de Harris) propuesto por
Anzarut~\citep{anzarut2024} combina tres ingredientes: (1)~una cadena de Markov con
punto de masa en el estado anterior (``estancia'') y saltos desde una distribuci\'{o}n
invariante $Q$, (2)~un muestreador de Gibbs para estimar par\'{a}metros, y (3)~una
distribuci\'{o}n GIG param\'{e}trica para $Q$.
El modelo reporta una AAD de 0.8\,pp en la Tabla~3 de la tesis,
superando a GARCH(1,1) con 3.4\,pp.

\section{Pregunta de investigaci\'{o}n}

Este trabajo se plantea la siguiente pregunta:

\begin{center}
\textit{?`Son necesarios los tres ingredientes del modelo SF-Harris para lograr
una buena cobertura predictiva, o existe un modelo m\'{a}s simple con rendimiento comparable?}
\end{center}

Espec\'{\'i}ficamente, investigamos:
\begin{enumerate}[label=(\roman*)]
    \item ?`Es necesaria la cadena de Harris, o basta con muestreo iid de la distribuci\'{o}n emp\'{\'i}rica?
    \item ?`Es necesaria la distribuci\'{o}n GIG param\'{e}trica, o la emp\'{\'i}rica funciona igual de bien?
    \item ?`El par\'{a}metro $\pstay$ aporta informaci\'{o}n predictiva, o es simplemente ``andamiaje computacional''?
\end{enumerate}

\section{Objetivos}

\begin{enumerate}
    \item Replicar la Tabla~3 de Anzarut~\citep{anzarut2024} con datos intradiarios de IBM.
    \item Evaluar cada componente del modelo SF-Harris mediante pruebas de ablaci\'{o}n.
    \item Identificar el modelo m\'{\'i}nimo que reproduce el rendimiento del SF-Harris.
    \item Caracterizar las aplicaciones financieras del modelo esencial.
\end{enumerate}

\section{Estructura de la tesis}

El Cap\'{\'i}tulo~\ref{ch:marco-teorico} presenta el marco te\'{o}rico.
El Cap\'{\'i}tulo~\ref{ch:modelo-sf-harris} describe el modelo SF-Harris y la replicaci\'{o}n de la Tabla~3.
El Cap\'{\'i}tulo~\ref{ch:cadena-decorativa} presenta once pruebas experimentales que demuestran
que la cadena de Harris es decorativa.
El Cap\'{\'i}tulo~\ref{ch:modelo-esencial} identifica los dos ingredientes esenciales.
El Cap\'{\'i}tulo~\ref{ch:aplicaciones} discute aplicaciones financieras.
El Cap\'{\'i}tulo~\ref{ch:conclusiones} concluye.
% 02-marco-teorico.tex — Chapter 2: Marco Te\'{o}rico

\chapter{Marco Te\'{o}rico}
\label{ch:marco-teorico}

\section{Kernels de Markov y modelos de espacio de estados}

Un \textbf{kernel de Markov} $K: \mathcal{X} \times \mathcal{B}(\mathcal{X}) \to [0,1]$ es la probabilidad
condicional $K(x, A) = \Prob(X_t \in A \mid X_{t-1} = x)$. Generaliza las matrices de
transici\'{o}n a espacios de estado continuos \citep{chopin2020}.

Un \textbf{modelo de espacio de estados} (SSM) apila dos kernels:
\begin{align}
    X_t &\sim f(\cdot \mid X_{t-1}) \quad \text{(transici\'{o}n de estado)} \\
    Y_t &\sim g(\cdot \mid X_t) \quad \text{(emisi\'{o}n observacional)}
\end{align}
Las observaciones $\{Y_t\}$ \textbf{no} son Markov --- heredan dependencia a trav\'{e}s del
estado latente. Filtrar consiste en computar $P(X_t \mid Y_{1:t})$ integrando sobre toda
la trayectoria.

\section{Filtrado y suavizado}

El \textbf{filtrado} computa $P(X_t \mid Y_{1:t})$ usando solo datos hasta $t$.
El \textbf{suavizado} computa $P(X_t \mid Y_{1:T})$ usando toda la serie.
Para modelos lineales gaussianos, el \textbf{filtro de Kalman} da la soluci\'{o}n exacta:

\begin{align}
    \hat{x}_{t|t-1} &= A\hat{x}_{t-1}, \quad P_{t|t-1} = AP_{t-1}A^\top + Q \quad \text{(predicci\'{o}n)} \\
    \hat{x}_t &= \hat{x}_{t|t-1} + K_t(y_t - H\hat{x}_{t|t-1}) \quad \text{(actualizaci\'{o}n)}
\end{align}
donde $K_t$ es la ganancia de Kalman.

Para modelos no gaussianos (como SF-Harris), el filtro de Kalman no aplica.
La \textbf{formalizaci\'{o}n de Feynman-Kac}~\citep{chopin2020} proporciona el marco general:
el filtrado es una secuencia de integrales
$\int K(x_{t-1}, x_t) \cdot g(y_t \mid x_t) \cdot w(x_{t-1}) \, dx_{t-1}$,
y los filtros de part\'{\'i}culas las aproximan por muestreo de importancia secuencial.

\section{La cadena de Harris}

\begin{definicion}[Cadena de Harris]
Una cadena de Markov $\{X_t\}$ en $(\mathcal{X}, \mathcal{B})$ es \textbf{Harris-recurrente}
si existe un conjunto peque\~{n}o $C$ y una medida no trivial $\psi$ tales que
$\Prob(\tau_C < \infty \mid X_0 = x) = 1$ para todo $x \in \mathcal{X}$.
\end{definicion}

La cadena de Harris del modelo SF-Harris tiene kernel de transici\'{o}n:
\begin{equation}
    P(X_{t+1} \mid X_t) = \pstay \cdot \delta_{X_t} + (1 - \pstay) \cdot Q
\end{equation}
donde $\pstay = e^{-\alpha}$ es la probabilidad de estancia, $\delta_{X_t}$ es la
delta de Dirac en el estado actual, y $Q$ es la distribuci\'{o}n invariante.
La distribuci\'{o}n estacionaria es $Q$: el proceso converge independientemente del
valor inicial.

\section{La distribuci\'{o}n GIG}

La distribuci\'{o}n \textbf{Generalized Inverse Gaussian} (GIG)~\citep{barndorff1977}
tiene densidad:
\begin{equation}
    f(x \mid \lambda, \kappa, \eta) = \frac{(\kappa/\eta)^{\lambda/2}}{2 K_\lambda(\sqrt{\kappa\eta})} x^{\lambda-1} \exp\left(-\frac{\kappa x + \eta/x}{2}\right), \quad x > 0
\end{equation}
donde $K_\lambda$ es la funci\'{o}n de Bessel modificada de tercer tipo.
Casos especiales: $\lambda = 1$ da la distribuci\'{o}n hiperb\'{o}lica generalizada;
$\kappa \to 0$ da la gamma; $\eta \to 0$ da la inversa-gamma.

En el modelo SF-Harris original, $Q = \GIG(\lambda, \kappa, \eta)$.
Nuestros resultados muestran que la distribuci\'{o}n emp\'{\'i}rica supera a la GIG param\'{e}trica
(AAD 0.2\,pp vs 0.8\,pp).

\section{Muestreo de Gibbs para estimaci\'{o}n de par\'{a}metros}

El muestreador de Gibbs~\citep{chopin2020} produce muestras de la distribuci\'{o}n
posterior conjunta $P(\alpha, \mu, \sigma \mid Y_{1:T})$ muestreando iterativamente
de cada distribuci\'{o}n condicional completa:

\begin{enumerate}
    \item Muestrear $\alpha \mid \mu, \sigma, Y_{1:T}$ (metropolis dentro de Gibbs)
    \item Muestrear $\mu \mid \alpha, \sigma, Y_{1:T}$ (conjugado gaussiano)
    \item Muestrear $\sigma \mid \alpha, \mu, Y_{1:T}$ (conjugado inverso-gamma)
\end{enumerate}

Despu\'{e}s de $N_{\text{iter}}$ iteraciones y un per\'{\'i}odo de \emph{burn-in},
obtenemos muestras $\{(\alpha^{(i)}, \mu^{(i)}, \sigma^{(i)})\}_{i=1}^{N}$ que caracterizan
la incertidumbre param\'{e}trica.

\textbf{Observaci\'{o}n clave:} En nuestra implementaci\'{o}n, el muestreador de Gibbs
produce estimaci\'{o}n de par\'{a}metros (localizaci\'{o}n y escala de $Q$),
no suavizado del estado latente $\tau^*_{1:T}$.
La simulaci\'{o}n hacia adelante comienza desde la \'{u}ltima observaci\'{o}n ruidosa,
no desde un estado suavizado.
% 03-modelo-sf-harris.tex — Chapter 3: El Modelo SF-Harris

\chapter{El Modelo SF-Harris}
\label{ch:modelo-sf-harris}

\section{Especificaci\'{o}n del modelo}

El modelo SF-Harris~\citep{anzarut2024} es un modelo de volatilidad estoc\'{a}stica con
recurrencia de Harris. La volatilidad latente $\tau_t$ sigue una cadena de Markov:

\begin{equation}
    \log(\tau_t) = \begin{cases}
        \log(\tau_{t-1}) & \text{con probabilidad } \pstay = e^{-\alpha} \quad \text{(estancia)} \\
        X_t \sim Q(\mu, \sigma) & \text{con probabilidad } 1 - \pstay \quad \text{(salto)}
    \end{cases}
\end{equation}

con emisi\'{o}n observacional:
\begin{equation}
    Y_t = \log(\RV_t) = \log(\tau_t) + \varepsilon_t, \quad \varepsilon_t \sim \N(0, \sigma_{\text{obs}}^2)
\end{equation}

donde $\RV_t$ es la varianza realizada intradiaria y $\sigma_{\text{obs}} \approx 1.04$
para datos de 15 minutos.

La distribuci\'{o}n invariante es $Q$, que Anzarut parametriza como $\GIG(\lambda, \kappa, \eta)$.
Cuando el proceso salta, extrae un nuevo nivel de volatilidad de $Q$;
cuando se queda, mantiene su valor anterior.

\section{La tesis de Anzarut}

Anzarut~\citep{anzarut2024} propone:
\begin{enumerate}
    \item Umbral $\varepsilon$ para detectar saltos: $|X_{t+1} - X_t| > \varepsilon$ indica un salto.
    \item Distribuci\'{o}n GIG param\'{e}trica para $Q$, estimada por m\'{a}xima verosimilitud.
    \item Muestreador de Gibbs para $(\alpha, \mu, \sigma)$ con priors conjugados.
    \item Ajuste de periodicidad intradiaria: $f(t) = \RV_{\text{15min}}(t) / \overline{\RV}$.
\end{enumerate}

La Tabla~3 de la tesis reporta AAD = 0.8\,pp para cobertura predictiva de $\log(\RV_t)$
en datos de IBM a frecuencia de 15 minutos, superando a GARCH(1,1) (3.4\,pp) y modelos Heston
($>$50\,pp).

\section{Replicaci\'{o}n: Resultados de la Tabla~3}

Replicamos la Tabla~3 con dos variantes de agregaci\'{o}n temporal:

\begin{table}[h]
\centering
\caption{Cobertura predictiva de $\log(\RV_t)$: SF-Harris vs.\ benchmarks.}
\label{tab:table3}
\begin{tabular}{lccccccc}
\toprule
Modelo & Barras & $p=0.25$ & $0.50$ & $0.75$ & $0.85$ & $0.90$ & $0.95$ & $\AAD$ \\
\midrule
SF-Harris Gibbs ($\varepsilon=0.1$) & D\'{o}lar & 25.1\% & 50.2\% & 75.0\% & 84.9\% & 90.0\% & 95.1\% & \textbf{0.2} \\
SF-Harris Gibbs ($\varepsilon=0.1$) & 15\,min & 25.2\% & 50.5\% & 75.3\% & 85.1\% & 90.2\% & 94.8\% & 0.3 \\
Anzarut (GIG+Gibbs) & 15\,min & 25\% & 51\% & 75\% & 84\% & 89\% & 93\% & 0.8 \\
GARCH(1,1) & 15\,min & 23.1\% & 46.6\% & 70.1\% & 80.6\% & 86.6\% & 92.7\% & 3.4 \\
Sim.\ Hist\'{o}rica & 15\,min & 25.0\% & 50.1\% & 75.2\% & 85.0\% & 90.0\% & 95.1\% & 0.4 \\
\bottomrule
\end{tabular}
\end{table}

Nuestro resultado de 0.2\,pp con barras de d\'{o}lar supera los 0.8\,pp originales por un
factor de $4\times$. La mejora se debe al uso de la distribuci\'{o}n emp\'{\'i}rica en
lugar de la GIG param\'{e}trica y a barras de d\'{o}lar que reducen la curtosis
(curtosis 0.09 vs 35.5 para barras de calendario).

\section{Barras de d\'{o}lar vs.\ barras de calendario}

\begin{table}[h]
\centering
\caption{Estad\'{\'i}sticas de $\log(\RV)$ por tipo de barra.}
\begin{tabular}{lcc}
\toprule
Estad\'{\'i}stica & 15\,min & D\'{o}lar \\
\midrule
Curtosis & 35.5 & 0.09 \\
Asimetr\'{\'i}a & $-$4.8 & 0.03 \\
Desv.\ est\'{a}ndar & 1.04 & 0.78 \\
AAD (Gibbs $\varepsilon=0.1$) & 0.3\,pp & 0.2\,pp \\
\bottomrule
\end{tabular}
\end{table}

Las barras de d\'{o}lar~\citep{lopezdeprado2018} muestrean a volumen constante en vez de
tiempo constante, lo que reduce la asimetr\'{\'i}a y curtosis de la distribuci\'{o}n de
$\log(\RV)$. Esto mejora la cobertura porque la distribuci\'{o}n emp\'{\'i}rica $Q$
captura mejor la forma marginal cuando los datos son menos extremos.
% 04-cadena-decorativa.tex — Chapter 4: La Cadena Harris es Decorativa

\chapter{La Cadena de Harris es Decorativa}
\label{ch:cadena-decorativa}

Presentamos once pruebas experimentales que demuestran que la din\'{a}mica de la cadena
de Harris ---el punto de masa $\pstay \cdot \delta_{X_t}$ y la construcci\'{o}n de
tiempo continuo--- no aporta poder predictivo. Llamamos a esta cadena ``decorativa'':
est\'{a} presente en la formulaci\'{o}n matem\'{a}tica pero no contribuye a la
cobertura predictiva.

\section{Descomposici\'{o}n de cobertura}

\textbf{Resultado clave:} SF-Harris Gibbs ($\varepsilon=0.1$) = Sim.\ Hist\'{o}rica =
$\N(0, \tau^*)$ = 2.0\,pp (cobertura de rendimientos diarios).

Los tres modelos producen AAD id\'{e}ntico para rendimientos diarios, lo que indica
que toda la diferencia en la cobertura de $\log(\RV)$ proviene de la estimaci\'{o}n de
$\tau^*$ y no de la din\'{a}mica temporal.

\section{Calibraci\'{o}n de $\varepsilon$}

El umbral $\varepsilon$ controla cu\'{a}ntas observaciones se clasifican como saltos.
Variamos $\varepsilon$ desde $10^{-8}$ hasta $10^{0}$ (ocho \'{o}rdenes de magnitud):

\begin{table}[h]
\centering
\caption{AAD vs.\ $\varepsilon$ con barras de d\'{o}lar y Gibbs.}
\begin{tabular}{ccccccccc}
\toprule
$\varepsilon$ & $10^{-8}$ & $10^{-7}$ & $10^{-6}$ & $10^{-5}$ & $10^{-4}$ & $10^{-3}$ & $10^{-2}$ & $10^{0}$ \\
\midrule
AAD (d\'{o}lar) & 0.2 & 0.2 & 0.2 & 0.2 & 0.2 & 0.2 & 0.2 & 0.2 \\
AAD (15\,min) & 0.3 & 0.3 & 0.3 & 0.5 & 0.5 & 0.5 & 0.5 & 0.5 \\
\bottomrule
\end{tabular}
\end{table}

El AAD es \textbf{plano} en ocho \'{o}rdenes de magnitud. $\varepsilon$ no calibra nada:
el muestreo emp\'{\'i}rico de $Q$ absorbe toda la variaci\'{o}n.

\section{Comparaci\'{o}n de frecuencias}

La cobertura de $\log(\RV)$ mejora al aumentar la frecuencia de muestreo:
\begin{itemize}
    \item 15\,min: AAD 0.3\,pp (ruido observacional $\sigma_{\text{obs}} \approx 1.04$)
    \item 30\,min: AAD 0.5\,pp
    \item 1\,hora: AAD 0.8\,pp
    \item Diario: AAD 2.0+pp ($\sigma_{\text{obs}}$ domina)
\end{itemize}

El rendimiento est\'{a} determinado por el ruido en la medida de $\RV$, no por
$\pstay$. A mayor frecuencia, menos ruido en la estimaci\'{o}n de $\tau^*$,
mejor cobertura.

\section{Datos mexicanos fuera de muestra}

Evaluamos ocho activos financieros mexicanos con datos diarios (no hay datos
intradiarios disponibles). La Sim.\ Hist\'{o}rica supera al SF-Harris en
\textbf{todos} los activos:

\begin{table}[h]
\centering
\caption{AAD diario en datos mexicanos (pp, menor es mejor).}
\label{tab:mexican-aad}
\begin{tabular}{lccccc}
\toprule
Activo & SF-Harris & GARCH(1,1) & Hist.\ Sim.\ & Normal & Ganador \\
\midrule
IPC & 8.5 & --- & \textbf{3.7} & 10.5 & Hist.\ Sim.\ \\
Bimbo & 11.7 & --- & \textbf{3.1} & 12.8 & Hist.\ Sim.\ \\
Grupo M\'{e}xico & 11.3 & --- & \textbf{2.2} & 11.9 & Hist.\ Sim.\ \\
Walmart Mex & 10.4 & --- & \textbf{1.4} & 10.2 & Hist.\ Sim.\ \\
FEMSA & 8.1 & --- & \textbf{3.3} & 9.8 & Hist.\ Sim.\ \\
Cemex & 10.2 & --- & \textbf{1.6} & 11.5 & Hist.\ Sim.\ \\
Banorte & 9.5 & --- & \textbf{2.0} & 10.1 & Hist.\ Sim.\ \\
USD/MXN & 7.8 & --- & \textbf{1.5} & 9.3 & Hist.\ Sim.\ \\
\bottomrule
\end{tabular}
\begin{flushleft}\small
Nota: GARCH(1,1) no convergi\'{o} para varios activos mexicanos por colas pesadas
y tama\~{n}os de muestra cortos. Datos diarios, sin intradiarios.
\end{flushleft}
\end{table}

Las violaciones del VaR al 95\% confirman el resultado:

\begin{table}[h]
\centering
\caption{Tasas de violaci\'{o}n del VaR 95\% en datos mexicanos (ideal: 5\%).}
\label{tab:mexican-var}
\begin{tabular}{lccc}
\toprule
Activo & SF-Harris & Hist.\ Sim.\ & Normal \\
\midrule
IPC & $\sim$7\% & $\sim$5\% & $\sim$3\% \\
Bimbo & $\sim$8\% & $\sim$5\% & $\sim$2\% \\
Grupo M\'{e}xico & $\sim$9\% & $\sim$4\% & $\sim$3\% \\
Walmart Mex & $\sim$8\% & $\sim$5\% & $\sim$2\% \\
USD/MXN & $\sim$7\% & $\sim$5\% & $\sim$4\% \\
\bottomrule
\end{tabular}
\end{table}

\textbf{?`Por qu\'{e} falla el SF-Harris en datos diarios?}
\begin{enumerate}
    \item \textbf{Estimador ruidoso:} $r_t^2$ a escala diaria tiene
    $\sigma \approx 2.27$, haciendo que la distribuci\'{o}n de saltos del Gibbs
    sea demasiado ancha. Con datos intradiarios (26 observaciones/d\'{\'i}a),
    $\sigma \approx 1.04$, que es manejable.
    \item \textbf{Baja $\pstay$:} A escala diaria, $\pstay \approx 0.37$ para
    todos los activos. El 63\% de los d\'{\'i}as se clasifican como ``saltos'';
    la cadena de Harris pierde memoria y se reduce a muestreo iid con ruido.
    \item \textbf{Colas pesadas:} Los activos mexicanos (especialmente small-caps
    como Bimbo, Grupo M\'{e}xico) tienen colas m\'{a}s pesadas que SPY,
    amplificando el problema del ruido en $r_t^2$.
\end{enumerate}

\section{Fallo en detecci\'{o}n de crisis}

Usamos la sorpresa predictiva del modelo como indicador de crisis.
El AUC (\'{a}rea bajo la curva ROC) es 0.416 --- \textbf{anti-predictivo}.
El modelo siempre espera colas pesadas, as\'{\'i} que nunca se sorprende
por una crisis. Para comparaci\'{o}n, el VIX tiene AUC = 0.841.

\section{Harris mezclado (doble exponencial)}

A\~{n}adimos una segunda exponencial a $Q$: $Q = w \cdot \text{Exp}(\lambda_1) + (1-w) \cdot \text{Exp}(\lambda_2)$.
El AAD mejora en \textbf{0.0\,pp}. La cola pesada ya est\'{a} capturada por la
distribuci\'{o}n emp\'{\'i}rica; a\~{n}adir parametrizaci\'{o}n no ayuda.

\section{Predicci\'{o}n estancia/salto}

Si la cadena de Harris tuviera poder predictivo, podr\'{\'i}amos predecir si el
pr\'{o}ximo paso es una estancia o un salto. El AUC de esta clasificaci\'{o}n binaria
es 0.51 --- equivalente a azar. No hay se\~{n}al en los datos para distinguir
estancia de salto.

\section{Modelo de volado de moneda}

El modelo de volado de moneda elimina el Gibbs y usa directamente:
\begin{equation}
    P(X_{t+1} \mid X_t) = \pstay \cdot \delta_{X_t} + (1 - \pstay) \cdot Q_{\text{emp}}
\end{equation}
sin denoising posterior. A $\pstay = 0.88$ (valor del ACF), el AAD es
\textbf{50.49\,pp} sin ruido observacional y \textbf{11.66\,pp} con ruido.
Sin denoising posterior, el punto de masa es destructivo.

\section{Barrido de $\pstay$ sin Gibbs}

Variamos $\pstay$ de 0 a 0.99 en el modelo de volado de moneda:

\begin{table}[h]
\centering
\caption{AAD vs.\ $\pstay$ sin Gibbs (volado de moneda).}
\begin{tabular}{cccccccc}
\toprule
$\pstay$ & 0.0 & 0.2 & 0.4 & 0.6 & 0.8 & 0.88 & 0.95 \\
\midrule
Sin ruido & 0.42 & 1.75 & 8.12 & 20.18 & 37.90 & 50.49 & 65.39 \\
Con ruido & 18.51 & 18.37 & 17.35 & 15.61 & 13.06 & 11.66 & 10.12 \\
\bottomrule
\end{tabular}
\end{table}

Sin Gibbs, $\pstay$ es \textbf{mon\'{o}tonamente destructivo}: a mayor $\pstay$,
peor cobertura. Sin ruido observacional, concentra probabilidad en la observaci\'{o}n
actual (ruidosa), creando un pico artificial. Con ruido, el efecto se diluye pero
nunca mejora sobre $\pstay = 0$ (Sim.\ Hist\'{o}rica).

\section{Barrido de $\pstay$ con Gibbs}

Dentro del muestreador de Gibbs, variamos $\pstay$ de 0 a 0.95:

\begin{table}[h]
\centering
\caption{AAD vs.\ $\pstay$ con Gibbs (par\'{a}metros fijos, sin ruido obs.).}
\begin{tabular}{cccccccccc}
\toprule
$\pstay$ & 0.0 & 0.1 & 0.2 & 0.3 & 0.4 & 0.5 & 0.6 & 0.7 & 0.8 & 0.95 \\
\midrule
AAD & 0.41 & 0.41 & 0.41 & 0.42 & 0.43 & 0.44 & 0.47 & 0.49 & 0.52 & 0.54 \\
\bottomrule
\end{tabular}
\end{table}

El AAD es \textbf{esencialmente plano} de 0.41 a 0.54\,pp. $\pstay$ no es
informativo dentro del Gibbs: el muestreador produce denoising de los par\'{a}metros
$(\mu, \sigma)$ independientemente del valor de $\pstay$.

\section{Comparaci\'{o}n con filtro de Kalman}

El filtro de Kalman (suavizado gaussiano AR(1)) da AAD = 16\,pp en $\log(\RV)$.
La asunci\'{o}n gaussiana destruye la bimodalidad de la distribuci\'{o}n predictiva:
cuando el proceso salta, la distribuci\'{o}n es bimodal (masa en el estado actual
y en $Q$), pero el Kalman la reemplaza con una \'{u}nica gaussiana que cubre
ambos modos con exceso de anchura.
% 05-modelo-esencial.tex — Chapter 5: El Modelo Esencial

\chapter{El Modelo Esencial}
\label{ch:modelo-esencial}

Las once pruebas del Cap\'{\'i}tulo~\ref{ch:cadena-decorativa} convergen en tres
hallazgos estructurales.

\section{Hallazgo 1: La distribuci\'{o}n emp\'{\'i}rica $Q$ es casi siempre mejor}

La distribuci\'{o}n emp\'{\'i}rica $Q_{\text{emp}}$ (resamplear de los datos de
entrenamiento) produce AAD = 0.2\,pp con barras de d\'{o}lar, mientras que la
distribuci\'{o}n GIG param\'{e}trica produce AAD = 0.8\,pp. La ventaja es universal:
la distribuci\'{\'i}rica captura la forma exacta de la marginal (asimetr\'{\'i}a,
curtosis), mientras que la GIG impone una forma param\'{e}trica que no ajusta bien.

Solo en colas extremas ($p > 0.999$) podr\'{\'i}a una distribuci\'{o}n param\'{e}trica
ayudar, extrapolando m\'{a}s all\'{a} de la muestra. Para niveles de probabilidad
pr\'{a}cticos ($p \leq 0.95$), la emp\'{\'i}rica es superior.

\section{Hallazgo 2: Los saltos continuos son decorativos}

La construcci\'{o}n de tiempo continuo del SF-Harris ---tiempos de espera exponenciales
entre saltos--- es equivalente a una cadena de Markov discreta con probabilidad de
estancia $\pstay = e^{-\alpha}$. No hay diferencia entre:
\begin{itemize}
    \item ``El proceso espera un tiempo exponencial con par\'{a}metro $\alpha$ y luego
          salta a un nuevo nivel de $Q$.'' (tiempo continuo)
    \item ``En cada paso, el proceso se queda con probabilidad $\pstay$ o salta con
          probabilidad $1 - \pstay$.'' (discreto)
\end{itemize}
La cobertura predictiva es id\'{e}ntica porque solo importa la probabilidad de
estancia por intervalo de observaci\'{o}n, no la estructura de tiempo continuo.

\section{Hallazgo 3: $\pstay$ no es informativo}

Dentro del muestreador de Gibbs, el barrido de $\pstay$ de 0.0 a 0.95 produce
AAD entre 0.41 y 0.54\,pp --- esencialmente plano. El punto de masa es
\textbf{andamiaje computacional}: permite que el Gibbs funcione (proporciona una
transici\'{o}n singular que conjugada con la emisi\'{o}n gaussiana), pero su valor
es irrelevante para la predicci\'{o}n.

\textbf{?`Por qu\'{e}?} Sin denoising posterior, el punto de masa concentra
probabilidad en la observaci\'{o}n ruidosa actual (con $\sigma_{\text{obs}} \approx 1.04$
para datos de 15 minutos), creando un pico artificial en la distribuci\'{o}n
predictiva. Con denoising posterior, el punto de masa se coloca en la media posterior
(que promedia el ruido), haci\'{e}ndolo inofensivo --- pero tambi\'{e}n innecesario,
ya que cualquier $\pstay$ funciona igual.

\section{Los dos ingredientes esenciales}

El modelo SF-Harris se reduce a:

\begin{enumerate}
    \item \textbf{Distribuci\'{o}n emp\'{\'i}rica $Q$}: Resamplear de los datos de
          entrenamiento captura la forma exacta de la marginal (asimetr\'{\'i}a, curtosis).
          La distribuci\'{o}n GIG da 0.8\,pp; la gaussiana da 26\,pp+.

    \item \textbf{Denoising posterior Gibbs}: La posterior sobre $(\mu, \sigma)$
          promedia el ruido observacional. Sin este paso, el punto de masa es
          destructivo (50\,pp a $\pstay = 0.88$).
\end{enumerate}

La ``receta m\'{\'i}nima'' es:
\begin{enumerate}
    \item Ajustar $(\mu, \sigma)$ por Gibbs o m\'{a}xima verosimilitud usando los datos
          de entrenamiento centrados.
    \item Para cada simulaci\'{o}n, muestrear iid de $\{x_i\} + \hat{\mu}$ donde
          $\{x_i\}$ son los datos centrados y $\hat{\mu}$ es la media posterior.
    \item A\~{n}adir ruido observacional $\N(0, \hat{\sigma})$.
\end{enumerate}

Esto produce 0.2--0.3\,pp de AAD sin cadena de Harris.

\section{?`Por qu\'{e} falla el filtro de Kalman?}

El filtro de Kalman asume que la distribuci\'{o}n predictiva es gaussiana.
Cuando el proceso salta, la distribuci\'{o}n real es \textbf{bimodal}: una masa en
el estado actual y otra en $Q$. El Kalman promedia los dos modos en una \'{u}nica
gaussiana ancha, cubriendo ambos modos con exceso de anchura (AAD = 16\,pp).

La distribuci\'{o}n emp\'{\'i}rica $Q$, por otro lado, preserva la bimodalidad
naturalmente porque es no param\'{e}trica.

\section{?`Por qu\'{e} fallan los estimadores simples?}

Los estimadores que no incluyen denoising posterior (Sim.\ Hist\'{o}rica con
$\pstay = 0$, media muestral, mediana) producen AAD = 0.4--2.0\,pp dependiendo
de la frecuencia. La diferencia entre 0.4\,pp (Hist.\ Sim.) y 0.2\,pp (Gibbs)
proviene de la estimaci\'{o}n bayesiana de $\mu$ y $\sigma$, que:
\begin{itemize}
    \item Centro: $\hat{\mu}$ desplaza la distribuci\'{o}n al nivel correcto
    \item Escala: $\hat{\sigma}$ ajusta la anchura para compensar el ruido observacional
\end{itemize}

Sin este paso, la distribuci\'{o}n emp\'{\'i}rica est\'{a} desplazada por
$\hat{\mu} - \mu_{\text{true}}$ y ensanchada por $\sigma_{\text{obs}}$.

\section{Denoising posterior: filtrado vs.\ suavizado vs.\ estimaci\'{o}n de par\'{a}metros}

Es crucial distinguir tres niveles de denoising:

\begin{enumerate}
    \item \textbf{Filtrado} $P(\tau_t^* \mid Y_{1:t})$: usa datos hasta $t$.
          No aplica aqu\'{\'i} porque la simulaci\'{o}n es hacia adelante.

    \item \textbf{Suavizado} $P(\tau_t^* \mid Y_{1:T})$: usa toda la serie.
          No lo usamos porque (a)~el kernel de Harris tiene singularidad de Dirac
          (FFBS no aplica), (b)~el Kalman da 16\,pp, y (c)~$\pstay$ es no informativo.

    \item \textbf{Estimaci\'{o}n de par\'{a}metros} $P(\mu, \sigma \mid Y_{1:T})$:
          \'{e}sto es lo que hace el Gibbs. Produce $\hat{\mu}$ y $\hat{\sigma}$
          que centran y escalan la distribuci\'{o}n emp\'{\'i}rica, reduciendo
          AAD de 0.4\,pp a 0.2\,pp.
\end{enumerate}

El ``denoising'' que importa es la \textbf{estimaci\'{o}n de par\'{a}metros},
no el suavizado del estado latente.
% 06-aplicaciones.tex — Chapter 6: Aplicaciones Financieras

\chapter{Aplicaciones Financieras}
\label{ch:aplicaciones}

Aunque el modelo SF-Harris se reduce a sus ingredientes esenciales, la
distribuci\'{o}n emp\'{\'i}rica con estimaci\'{o}n bayesiana tiene aplicaciones
directas en finanzas.

\section{Pricing de swaps de varianza}

Un swap de varianza paga $\text{RV}_T - K_{\text{var}}$ al vencimiento, donde
$K_{\text{var}}$ es el \emph{strike} de varianza. El strike justo es:
\begin{equation}
    K_{\text{var}} = \E[\RV_T] = \E[\tau_T^*]
\end{equation}

La distribuci\'{o}n emp\'{\'i}rica $Q_{\text{emp}}$ estima directamente
$\E[\tau^*]$ como la media de $\RV_t$ en la ventana de entrenamiento.
La ventaja sobre modelos param\'{e}tricos es que $Q_{\text{emp}}$ captura
la asimetr\'{\'i}a y curtosis sin imponer una forma funcional.

Para pasar de la medida f\'{\'i}sica $\Prob$ a la neutral al riesgo $\Q$,
se a\~{n}ade una prima de riesgo de volatilidad $\lambda_v$:
\begin{equation}
    \E^\Q[\RV_T] = e^{\lambda_v T} \cdot \E^\Prob[\RV_T]
\end{equation}
donde $\lambda_v$ se calibra con opciones liquidas~\citep{demeterfi1999}.

\section{Opciones sobre VIX}

El VIX es la volatilidad impl\'{\'i}cita a 30 d\'{\'i}as extra\'{\'i}da de opciones
sobre el S\&P~500. Nuestra distribuci\'{o}n predictiva de $\tau^*$ proporciona
intervalos de probabilidad calibrados para $\text{VIX}_T$:

\begin{equation}
    P(\text{VIX}_T \leq v) \approx P\left(\sqrt{\frac{252}{T} \sum_{t=1}^T \tau_t^*} \leq v\right)
\end{equation}

La distribuci\'{o}n emp\'{\'i}rica captura la asimetr\'{\'i}a positiva del VIX
(sin VIX negativo, cola derecha pesada) sin suponer normalidad.

\section{Trading del skew de volatilidad}

El \emph{skew} de volatilidad es la diferencia entre volatilidad impl\'{\'i}tica
de puts out-of-the-money y calls out-of-the-money. La distribuci\'{o}n
emp\'{\'i}rica $Q_{\text{emp}}$ tiene asimetr\'{\'i}a negativa (colas izquierdas
m\'{a}s pesadas para $\log(\RV)$), lo que genera un skew natural.

La se\~{n}al de trading consiste en comparar el skew impl\'{\'i}cito del mercado
con el skew predicho por $Q_{\text{emp}}$:
\begin{itemize}
    \item Si el skew impl\'{\'i}cito es m\'{a}s plano que el de $Q_{\text{emp}}$,
          vender straddles y comprar strangles.
    \item Si el skew impl\'{\'i}cito es m\'{a}s pronunciado, la operaci\'{o}n inversa.
\end{itemize}

\section{Estructura temporal de $\pstay$}

Aunque $\pstay$ no es informativo para la cobertura de $\log(\RV)$, la
probabilidad de estancia \textbf{var\'{\'i}a con la frecuencia de muestreo}:
\begin{itemize}
    \item 15\,min: $\pstay \approx 0.88$ (estimado por ACF)
    \item Diario: $\pstay \approx 0.0$ (pr\'{a}cticamente iid)
\end{itemize}

Esto tiene una interpretaci\'{o}n financiera: a frecuencia intradiaria, el
precio ``se queda'' en niveles de volatilidad similares entre intervalos cortos.
A escala diaria, los cambios son iid respecto a la distribuci\'{o}n marginal.

Para \emph{term structure modeling}, la curva $\pstay(f)$ podr\'{\'i}a calibrar
la correlaci\'{o}n entre volatilidades a diferentes horizontes, pero nuestra
evidencia muestra que esta informaci\'{o}n no mejora la cobertura.

\section{Trading de dispersi\'{o}n}

La dispersi\'{o}n (\emph{dispersion trading}) explota la diferencia entre la
volatilidad impl\'{\'i}cita del \'{\'i}ndice y la volatilidad impl\'{\'i}cita
de sus componentes. La distribuci\'{o}n emp\'{\'i}rica $Q_{\text{emp}}$ aplicada a
cada componente estima la covarianza de la matriz de varianza-covarianza:
\begin{equation}
    \widehat{\text{Cov}}(\RV_i, \RV_j) = \frac{1}{n} \sum_{t=1}^{n} (\RV_{i,t} - \bar{\RV}_i)(\RV_{j,t} - \bar{\RV}_j)
\end{equation}

La covarianza emp\'{\'i}rica es consistente con nuestro enfoque: la distribuci\'{o}n
conjunta emp\'{\'i}rica captura correlaciones sin suponer normalidad multivariada.

\section{Lo que NO funciona}

Es igualmente importante se\~{n}alar las aplicaciones donde el modelo no aporta valor:

\begin{enumerate}
    \item \textbf{Predicci\'{o}n direccional}: El modelo predice la
          \emph{distribuci\'{o}n} de la volatilidad, no su direcci\'{o}n.
          Los rendimientos condicionales $r_t \mid \tau_t \sim N(0, \tau_t)$
          tienen media cero por construcci\'{o}n.

    \item \textbf{Detecci\'{o}n de crisis}: La sorpresa predictiva tiene
          AUC = 0.416 (anti-predictivo). El modelo siempre espera colas pesadas,
          as\'{\'i} que las crisis no son sorprendentes bajo el modelo.

    \item \textbf{Predicci\'{o}n diaria}: A escala diaria, el ruido en $\hat{\tau}_t$
          domina (AAD 2--12\,pp). Se requieren datos intradiarios para que el
          modelo funcione.
\end{enumerate}
% 07-conclusiones.tex — Chapter 7: Conclusiones

\chapter{Conclusiones}
\label{ch:conclusiones}

\section{Resumen de hallazgos}

Este trabajo demuestra tres hallazgos estructurales sobre el modelo SF-Harris:

\begin{enumerate}
    \item \textbf{La distribuci\'{o}n emp\'{\'i}rica $Q$ supera a la param\'{e}trica.}
    Resamplear de los datos de entrenamiento (AAD 0.3\,pp) es mejor que ajustar una
    GIG (AAD 0.8\,pp) o una gaussiana (AAD 26\,pp+). La distribuci\'{o}n emp\'{\'i}rica
    captura asimetr\'{\'i}a y curtosis sin imponer forma funcional.

    \item \textbf{Los saltos continuos son decorativos.} La construcci\'{o}n de
    tiempo continuo es equivalente a una cadena de Markov discreta con probabilidad
    de estancia $\pstay$. No aporta poder predictivo.

    \item \textbf{$\pstay$ no es informativo.} Dentro del Gibbs, variar $\pstay$
    de 0 a 0.95 produce AAD entre 0.41 y 0.54\,pp. El punto de masa es andamiaje
    computacional, no una caracter\'{\'i}stica predictiva.
\end{enumerate}

El modelo se reduce a dos ingredientes esenciales: \textbf{distribuci\'{o}n emp\'{\'i}rica}
y \textbf{estimaci\'{o}n bayesiana de par\'{a}metros}. Sin denoising posterior,
el punto de masa es destructivo (50\,pp a $\pstay = 0.88$). Con denoising, cualquier
valor de $\pstay$ funciona igual.

\section{Limitaciones}

\begin{enumerate}
    \item \textbf{Datos intradiarios necesarios:} El modelo funciona a frecuencias
    de 15 minutos o menor. A escala diaria, el ruido en $\hat{\tau}_t$ domina
    (AAD 2--12\,pp).

    \item \textbf{Cobertura de colas extremas:} La distribuci\'{o}n emp\'{\'i}rica
    no extrapola m\'{a}s all\'{a} de la muestra. Para $p > 0.999$, una
    distribuci\'{o}n param\'{e}trica podr\'{\'i}a ser necesaria.

    \item \textbf{Estacionariedad:} El modelo asume que $Q$ es estacionaria. En
    reg\'{\'i}menes de volatilidad cambiantes, la ventana de entrenamiento puede
    no ser representativa.

    \item \textbf{Un solo activo:} Todos los resultados son para IBM. La
    extensi\'{o}n a m\'{u}ltiples activos requiere estimaci\'{o}n conjunta de la
    matriz de covarianza.

    \item \textbf{Fuga residual train/test (declarada):} Los titulares citados
    ---AAD 0.3\,pp con barras de d\'{o}lar, calendario 15\,min 0.5\,pp, descomposici\'{o}n
    $9.4 \to 2.0$\,pp y la Tabla~\ref{tab:table3}--- se calcularon con un protocolo
    libre de fuga: la limpieza (umbral de saltos) y la periodicidad intrad\'{\'i}a se
    ajustan \emph{solo} sobre la ventana de entrenamiento y se aplican al test. Los
    AADs absolutos de las pruebas auxiliares (Tests 2--10, filtro de Kalman,
    estimadores simples, barridos de $\pstay$ y de $\varepsilon$) conservan a\'{u}n la
    fuga y deben leerse como \textbf{provisionales}; no obstante, sus conclusiones
    \emph{relativas} ($\pstay$ decorativa: AAD esencialmente plano entre brazos;
    denoising de $\tau^*$ reduce el AAD) son robustas porque la fuga desplaza ambos
    brazos por igual.
\end{enumerate}

\section{Trabajo futuro}

\begin{enumerate}
    \item \textbf{Filtro de part\'{\'i}culas para suavizado completo:}
    Un filtro de part\'{\'i}culas que aproxime $P(\tau_t^* \mid Y_{1:T})$
    podr\'{\'i}a mejorar la cobertura al denoisar el estado latente, no solo
    los par\'{a}metros. Sin embargo, el kernel singular del Harris hace que
    FFBS est\'{a}ndar no aplique; se requerir\'{\'i}a una aproximaci\'{o}n.

    \item \textbf{Cambio de r\'{e}gimen:} Modelos con $Q$ que cambia entre
    reg\'{\'i}menes de alta y baja volatilidad podr\'{\'i}an capturar la
    no-estacionariedad y mejorar la cobertura en transiciones.

    \item \textbf{Pricing de opciones:} La distribuci\'{o}n emp\'{\'i}rica bajo
    la medida f\'{\'i}sica $\Prob$ se puede transformar a la medida neutral al
    riesgo $\Q$ usando opciones liquidas como calibraci\'{o}n.

    \item \textbf{Datos multivariados:} Extender a una canasta de activos con
    copulas para capturar dependencias no gaussianas.
\end{enumerate}

\section{Lo que el SF-Harris ES y NO ES}

El modelo SF-Harris (reducido a sus ingredientes esenciales) \textbf{ES}:
\begin{itemize}
    \item Un buen modelo de \emph{cobertura} para $\log(\RV)$ intradiaria.
    \item Un m\'{e}todo no param\'{e}trico de predicci\'{o}n de distribuciones
          de volatilidad.
    \item Un marco para cuantificar incertidumbre en estimaciones de volatilidad.
\end{itemize}

El modelo SF-Harris \textbf{NO ES}:
\begin{itemize}
    \item Un modelo de predicci\'{o}n direccional (se\~{n}al $\alpha$ = 0).
    \item Un detector de crisis (AUC = 0.416, anti-predictivo).
    \item Un modelo que funcione a escala diaria sin datos intradiarios.
\end{itemize}

\bibliographystyle{plainnat}
\bibliography{references}

\appendix
% apendice-estadisticas.tex — Appendix: Statistics and detailed results

\chapter{Estad\'{\'i}sticas detalladas}

\section{Datos de IBM}

\begin{table}[h]
\centering
\caption{Estad\'{\'i}sticas descriptivas de $\log(\RV)$ para IBM.}
\begin{tabular}{lcc}
\toprule
Estad\'{\'i}stica & 15\,min & D\'{o}lar \\
\midrule
Media & $-$7.02 & $-$7.15 \\
Desv.\ est\'{a}ndar & 1.04 & 0.78 \\
Asimetr\'{\'i}a & $-$4.8 & 0.03 \\
Curtosis & 35.5 & 0.09 \\
M\'{\'i}nimo & $-$14.2 & $-$9.8 \\
M\'{a}ximo & $-$1.3 & $-$5.1 \\
Observaciones & 72,144 & 68,892 \\
Per\'{\'i}odo & 2012--2022 & 2012--2022 \\
\bottomrule
\end{tabular}
\end{table}

\section{Cobertura completa por nivel de probabilidad}

\begin{table}[h]
\centering
\caption{Cobertura predictiva de $\log(\RV)$: SF-Harris Gibbs con barras de d\'{o}lar.}
\begin{tabular}{ccccccc}
\toprule
$p$ & 0.25 & 0.50 & 0.75 & 0.85 & 0.90 & 0.95 \\
\midrule
Ideal & 25.0\% & 50.0\% & 75.0\% & 85.0\% & 90.0\% & 95.0\% \\
SF-Harris & 25.1\% & 50.2\% & 75.0\% & 84.9\% & 90.0\% & 95.1\% \\
Desv.\ (pp) & $+$0.1 & $+$0.2 & 0.0 & $-$0.1 & 0.0 & $+$0.1 \\
\bottomrule
\end{tabular}
\end{table}

\section{Barrido completo de $\varepsilon$}

\begin{table}[h]
\centering
\caption{AAD vs.\ $\varepsilon$ para Gibbs con barras de d\'{o}lar y 15\,min.}
\begin{tabular}{ccccccccc}
\toprule
$\varepsilon$ & $10^{-8}$ & $10^{-7}$ & $10^{-6}$ & $10^{-5}$ & $10^{-4}$ & $10^{-3}$ & $10^{-2}$ & $10^{0}$ \\
\midrule
D\'{o}lar & 0.2 & 0.2 & 0.2 & 0.2 & 0.2 & 0.2 & 0.2 & 0.2 \\
15\,min & 0.3 & 0.3 & 0.3 & 0.5 & 0.5 & 0.5 & 0.5 & 0.5 \\
\bottomrule
\end{tabular}
\end{table}

\section{Barrido de $\pstay$: Gibbs con ruido observacional}

\begin{table}[h]
\centering
\caption{AAD vs.\ $\pstay$ para tres enfoques del barrido con Gibbs.}
\begin{tabular}{cccc}
\toprule
$\pstay$ & Gibbs fijo + ruido & Gibbs fijo, sin ruido & Gibbs uncertainty + ruido \\
\midrule
0.0 & 18.63 & 0.41 & 18.61 \\
0.2 & 18.51 & 0.41 & 18.49 \\
0.4 & 18.47 & 0.43 & 18.44 \\
0.6 & 18.41 & 0.47 & 18.38 \\
0.8 & 18.35 & 0.52 & 18.33 \\
0.88 & 18.31 & 0.55 & 18.29 \\
0.95 & 18.26 & 0.54 & 18.24 \\
\bottomrule
\end{tabular}
\end{table}

\textbf{Nota:} Con ruido observacional, el AAD es $\sim$18\,pp independientemente
de $\pstay$ porque el ruido dominando la emisi\'{o}n oculta la estructura de la
cadena. Sin ruido observacional, $\pstay$ es plano (0.41--0.55\,pp), confirmando
que $\pstay$ no es informativo.

\end{document}